% =====================================================================
\chapter{Coste del gradiente y productos vector-jacobiano}
\label{tema:4}
% =====================================================================

Este tema responde a la pregunta que justifica el módulo: por qué el modo inverso calcula el gradiente de una función de millones de variables con un coste comparable al de evaluarla. Se demuestra el principio del gradiente barato, se compara el coste de los dos modos mediante la asociatividad del producto de jacobianos, se formulan las reglas locales de los sistemas reales como productos vector-jacobiano y se deduce de ellas la regla $C\approx 6ND$ del cómputo de entrenamiento. Se siguen el capítulo~4 de~\cite{gw}, la §3.2 de~\cite{baydin}, la lección~6 de~\cite{grosse}, el apartado «Gradients for vectorized operations» de~\cite{cs231n} y el artículo de Kaplan \emph{et al.}~\cite{kaplan}.

\section{El principio del gradiente barato}

\begin{teorema}[Gradiente barato]\label{teo:barato}
Sea $f\colon\R^{n}\to\R$ dada por una traza con operaciones del conjunto $\Phi$ de la \cref{def:elementales}, cada una con a lo sumo $k$ argumentos. Existe una constante $\omega$, que depende solo de $\Phi$ y no de $n$ ni de la traza, tal que el modo inverso calcula $\nabla f(x)$ con a lo sumo $\omega\cdot\ops(f)$ operaciones.
\end{teorema}

\begin{proof}
Sea $\ell$ el número de operaciones de la traza y $E$ el de aristas del grafo. Cada operación cuesta al menos una unidad, luego $\ops(f)\ge\ell$, y tiene a lo sumo $k$ argumentos, luego $E\le k\ell$.
\begin{enumerate}[label=\arabic*.]
  \item La pasada hacia delante cuesta $\ops(f)$.
  \item En la pasada hacia atrás, la operación $i$, con $r_i\le k$ argumentos, evalúa $r_i$ derivadas locales sobre valores guardados, cada una de coste a lo sumo $c_\Phi$ (\cref{def:elementales}), y ejecuta $r_i$ productos y $r_i$ sumas. Su coste es a lo sumo $(c_\Phi+2)\,r_i$.
  \item Sumando sobre las operaciones, la pasada hacia atrás cuesta a lo sumo $(c_\Phi+2)\sum_ir_i=(c_\Phi+2)E\le(c_\Phi+2)k\ell\le(c_\Phi+2)k\,\ops(f)$.
\end{enumerate}
El total es a lo sumo $\bigl(1+(c_\Phi+2)k\bigr)\ops(f)$, y basta tomar $\omega=1+(c_\Phi+2)k$. La corrección del resultado es el \cref{teo:inverso}.
\end{proof}

La constante de la demostración es pesimista. Un análisis más fino de cada operación da cotas pequeñas: Baydin \emph{et al.}, siguiendo a Griewank y Walther, indican que el cociente entre el coste del modo inverso y el de la función está por debajo de~6 y suele estar entre~2 y~3~\cite[§3.2]{baydin}, \cite[cap.~4]{gw}. Para funciones racionales, un resultado de este tipo fue demostrado por Baur y Strassen~\cite{baurstrassen}. Lo esencial es que la cota no depende de $n$: diez derivadas y $10^{11}$ derivadas cuestan, en relación con $f$, lo mismo.

\begin{corolario}[Elección del modo]\label{cor:eleccion}
Para $f\colon\R^{n}\to\R^{m}$, el jacobiano completo cuesta del orden de $n\cdot\ops(f)$ en modo directo y de $m\cdot\ops(f)$ en modo inverso. Conviene el directo si $n\ll m$ y el inverso si $m\ll n$. Entrenar un LLM corresponde a $n\approx 10^{11}$ y $m=1$.
\end{corolario}

\begin{proof}
Es la combinación del \cref{cor:coste-directo} y la \cref{prop:coste-dual} con la \cref{prop:vjp} y el \cref{teo:barato}, aplicado este a cada una de las $m$ funciones escalares $\bar y\T f$ con $\bar y=e_r$.
\end{proof}

El precio del modo inverso es la memoria: la pasada hacia atrás necesita los valores de la pasada hacia delante. En el modo directo, en cambio, cada tangente se descarta en cuanto deja de usarse. El \cref{tema:5} estudia este compromiso.

\section{Asociatividad del producto de jacobianos}

Sea $f=f_L\circ\dots\circ f_1$ una composición de funciones diferenciables. Por la regla de la cadena (módulo~4), $J_f=J_{f_L}\cdots J_{f_1}$, con cada jacobiano evaluado en el punto correspondiente. El producto de matrices es asociativo y el resultado no depende del orden en que se agrupen los factores, pero el coste sí.

\begin{proposicion}[Coste de las dos agrupaciones]\label{prop:asociatividad}
Sea $f\colon\R^{m}\to\R$ con $J_f=J_L J_{L-1}\cdots J_1$, donde $J_L\in\R^{1\times m}$ y $J_1,\dots,J_{L-1}\in\R^{m\times m}$, con $L\ge 2$, y cuéntese como coste de un producto $(p\times q)(q\times r)$ el número $pqr$ de multiplicaciones escalares. Entonces:
\begin{enumerate}
  \item agrupando de derecha a izquierda, $J_L\bigl(J_{L-1}(\cdots(J_2J_1))\bigr)$, el coste es $(L-2)m^{3}+m^{2}$;
  \item agrupando de izquierda a derecha, $\bigl(((J_LJ_{L-1})J_{L-2})\cdots\bigr)J_1$, el coste es $(L-1)m^{2}$.
\end{enumerate}
\end{proposicion}

\begin{proof}
(i) Los $L-2$ productos $J_2J_1$, $J_3(J_2J_1)$, \dots, $J_{L-1}(\cdots)$ son de matrices $m\times m$ y cuestan $m^{3}$ cada uno; el último, $J_L$ por una matriz $m\times m$, cuesta $1\cdot m\cdot m$. (ii) Cada producto parcial es un vector fila $1\times m$, y cada uno de los $L-1$ productos por una matriz $m\times m$ cuesta $m^{2}$.
\end{proof}

La agrupación (i) es la del modo directo con todas las semillas a la vez: propaga la matriz $\partial v/\partial x$, de tamaño $m\times m$. La agrupación (ii) es la del modo inverso: lo que viaja hacia atrás es siempre un vector fila, el adjunto. Con $L=3$ y $m=100$, el cociente de costes es $(100^{3}+100^{2})/(2\cdot 100^{2})=50{,}5$, y en general crece linealmente con $m$. Para composiciones que no son cadenas, encontrar la agrupación óptima es un problema difícil, estudiado en~\cite[cap.~9 y~10]{gw}.

\section{Productos vector-jacobiano}

Los sistemas reales no trabajan con escalares, sino con operaciones sobre tensores, y nunca construyen sus jacobianos~\cite{grosse}.

\begin{definicion}[Producto vector-jacobiano]\label{def:vjp}
Sea $g\colon\R^{p}\to\R^{q}$ una operación con entrada $u$ y salida $w=g(u)$, y sea $\bar w\in\R^{1\times q}$ el adjunto de la salida. El \emph{producto vector-jacobiano} (VJP) de $g$ en $u$ es la aplicación lineal $\bar w\mapsto\bar u=\bar w\,J_g(u)\in\R^{1\times p}$.
\end{definicion}

Por el \cref{teo:inverso}, aplicado por bloques, el modo inverso sobre un grafo de operaciones tensoriales solo necesita, para cada tipo de operación, una función que calcule su VJP. Si la entrada o la salida son matrices, se aplanan para la definición y el resultado se devuelve con la forma original.

\begin{ejemplo}[Operaciones componente a componente]\label{ej:vjp-componente}
Si $w=\phi(u)$ se aplica componente a componente, $J_g(u)=\operatorname{diag}(\phi'(u))$ y el VJP es $\bar u=\bar w\odot\phi'(u)$, con coste $O(p)$ en lugar de los $O(p^{2})$ que exigiría multiplicar por la matriz diagonal~\cite{grosse}. Para $\phi=\exp$, $\bar u=\bar w\odot w$: basta haber guardado la salida.
\end{ejemplo}

\begin{proposicion}[VJP de la capa lineal]\label{prop:vjp-lineal}
Sea $Y=XW$ con $X\in\R^{B\times n}$, $W\in\R^{n\times m}$, y sea $\bar Y\in\R^{B\times m}$ el adjunto de la salida. Entonces
\[
\bar X=\bar Y\,W\T\in\R^{B\times n},\qquad \bar W=X\T\,\bar Y\in\R^{n\times m}.
\]
\end{proposicion}

\begin{proof}
Se tiene $Y_{rs}=\sum_iX_{ri}W_{is}$, luego $\partial Y_{rs}/\partial X_{ri}=W_{is}$ y $\partial Y_{rs}/\partial W_{is}=X_{ri}$. Por la regla de la cadena, $\bar X_{ri}=\sum_s\bar Y_{rs}W_{is}=(\bar YW\T)_{ri}$ y $\bar W_{is}=\sum_r\bar Y_{rs}X_{ri}=(X\T\bar Y)_{is}$. Es el resultado obtenido en el módulo~4 (tema~3) y en el módulo~9 (tema~5), ahora como regla local de un sistema de AD.
\end{proof}

\begin{ejemplo}\label{ej:vjp-lineal}
Sean $X=(1,2)$, $W=\begin{psmallmatrix}1&0\\-1&1\end{psmallmatrix}$ y $\bar Y=(1,-1)$. Entonces $Y=(-1,2)$, $\bar X=\bar YW\T=(1,-2)$ y $\bar W=X\T\bar Y=\begin{psmallmatrix}1&-1\\2&-2\end{psmallmatrix}$. Ambos gradientes tienen la forma de la variable a la que corresponden.
\end{ejemplo}

\begin{proposicion}[VJP de la softmax]\label{prop:vjp-softmax}
Sea $p=\softmax(z)$ con $z\in\R^{1\times K}$ y sea $\bar p\in\R^{1\times K}$. Entonces
\[
\bar z=\bar p\,J_{\softmax}(z)=p\odot\bigl(\bar p-\langle\bar p,p\rangle\,\mathbf{1}\bigr),
\]
donde $\langle\bar p,p\rangle=\sum_k\bar p_kp_k$ y $\mathbf 1=(1,\dots,1)$. El cálculo cuesta $O(K)$ operaciones.
\end{proposicion}

\begin{proof}
Como se demostró en el módulo~4 (tema~3), $\partial p_i/\partial z_j=p_i(\ind[i=j]-p_j)$. Luego $\bar z_j=\sum_i\bar p_i\,p_i(\ind[i=j]-p_j)=\bar p_jp_j-p_j\sum_i\bar p_ip_i=p_j\bigl(\bar p_j-\langle\bar p,p\rangle\bigr)$. El producto escalar cuesta $O(K)$ y se calcula una vez; cada componente cuesta después $O(1)$.
\end{proof}

\begin{corolario}[Softmax con entropía cruzada]\label{cor:ce}
Si $L=-\ln p_y$, entonces $\bar z=p-e_y$.
\end{corolario}

\begin{proof}
El VJP de $p\mapsto-\ln p_y$ con $\bar L=1$ es $\bar p=-e_y/p_y$, y $\langle\bar p,p\rangle=-1$. Por la \cref{prop:vjp-softmax}, $\bar z_j=p_j(\bar p_j+1)$, que vale $p_y(-1/p_y+1)=p_y-1$ si $j=y$ y $p_j$ si $j\ne y$.
\end{proof}

En la práctica, la softmax y la entropía cruzada se tratan como una sola operación con VJP $p-e_y$, que además es numéricamente estable: es lo que hace \texttt{F.cross\_entropy} en PyTorch, que recibe logits y no probabilidades (\cref{tema:6}).

\section{El cómputo de un paso de entrenamiento}

Se cuenta como FLOP una suma o un producto en coma flotante. Un producto de matrices $(p\times q)(q\times r)$ cuesta $pr(2q-1)\approx 2pqr$ FLOPs.

\begin{proposicion}[Coste de una capa lineal]\label{prop:flops}
Para $Y=XW$ con $X\in\R^{B\times n}$ y $W\in\R^{n\times m}$, la pasada hacia delante cuesta unos $2Bnm$ FLOPs y la pasada hacia atrás unos $4Bnm$: $2Bnm$ para $\bar X$ y $2Bnm$ para $\bar W$. El total es $6Bnm=6\cdot(\text{parámetros})\cdot(\text{filas del lote})$.
\end{proposicion}

\begin{proof}
$XW$ es $(B\times n)(n\times m)$; $\bar YW\T$ es $(B\times m)(m\times n)$; $X\T\bar Y$ es $(n\times B)(B\times m)$. Los tres productos tienen el mismo número de multiplicaciones, $Bnm$, y cuestan unos $2Bnm$ FLOPs cada uno. La capa tiene $nm$ pesos.
\end{proof}

\begin{corolario}[Regla $6ND$]\label{cor:6nd}
Si el cómputo de un modelo con $N$ parámetros está dominado por productos de matrices en los que cada parámetro actúa una vez por token, entrenarlo con $D$ tokens cuesta $C\approx 6ND$ FLOPs: $2ND$ en las pasadas hacia delante y $4ND$ en las pasadas hacia atrás.
\end{corolario}

\begin{proof}
Se aplica la \cref{prop:flops} a cada capa, con $B$ igual al número de tokens procesados, y se suma sobre las capas y sobre los pasos.
\end{proof}

La aproximación ignora las operaciones que no son productos por pesos, como el producto de la atención entre consultas y claves, que crece con la longitud del contexto, y el hecho de que la primera capa no necesita $\bar X$. Kaplan \emph{et al.} la usan como base de las leyes de escalado~\cite{kaplan}, que se estudian en el módulo~16. Por ejemplo, un modelo de $N=10^{8}$ parámetros entrenado con $D=2\cdot 10^{9}$ tokens requiere $C\approx 6\cdot 10^{8}\cdot 2\cdot 10^{9}=1{,}2\cdot 10^{18}$ FLOPs; a $10^{14}$ FLOP/s efectivos, unos $1{,}2\cdot 10^{4}$~s, algo más de tres horas.

\begin{lecturas}
\item Griewank y Walther~\cite{gw}, cap.~4 «Memory Issues and Complexity Bounds».
\item Baydin \emph{et al.}~\cite{baydin}, §3.2, discusión final sobre el coste de los dos modos.
\item Grosse~\cite{grosse}, lección~6, apartado sobre productos vector-jacobiano.
\item Stanford CS231n~\cite{cs231n}, «Backpropagation, Intuitions», apartado «Gradients for vectorized operations».
\item Kaplan \emph{et al.}~\cite{kaplan}, recuento de parámetros y de cómputo de un Transformer y aproximación $C\approx 6ND$.
\end{lecturas}

\section*{Ejercicios del tema 4}
\addcontentsline{toc}{section}{Ejercicios del tema 4}

\begin{ejercicio}{1}{eje:4-modo}
Se dispone de $f\colon\R^{1000}\to\R$ y de $g\colon\R\to\R^{1000}$. ¿Qué modo de AD usaría para el jacobiano de cada una y cuántas pasadas necesitaría? ¿Cuántas necesitaría con el otro modo?
\end{ejercicio}

\begin{ejercicio}{1}{eje:4-d2l}
¿Cuándo conviene el modo directo y cuándo el inverso? Responda atendiendo a tres criterios: la cantidad de datos intermedios que hay que guardar, la posibilidad de ejecutar pasos en paralelo y el tamaño de las matrices y vectores que intervienen.
\fuente{adaptado de Zhang \emph{et al.}~\cite{d2l}, §2.5 «Automatic Differentiation», ejercicio~8.}
\end{ejercicio}

\begin{ejercicio}{2}{eje:4-cadena}
Sea $\nabla f\T=c\,B\,A$ con $c\in\R^{1\times 100}$ y $A,B\in\R^{100\times 100}$. Cuente las multiplicaciones escalares de las dos agrupaciones posibles y su cociente. Generalice a $c\in\R^{1\times m}$ y $A,B\in\R^{m\times m}$.
\end{ejercicio}

\begin{ejercicio}{2}{eje:4-lineal}
Para la capa $Y=XW$ con $X=(2,-1)$, $W=\begin{psmallmatrix}1&2&0\\0&1&-1\end{psmallmatrix}$ y adjunto de la salida $\bar Y=(1,0,2)$, calcule $Y$, $\bar X$ y $\bar W$, y compruebe las formas. Verifique una componente de $\bar W$ derivando $L=\bar Y\,Y\T$ directamente.
\end{ejercicio}

\begin{ejercicio}{2}{eje:4-cs336}
Un modelo de $N=10^{9}$ parámetros se entrena durante $10^{5}$ pasos con lotes de 512 secuencias de 1024 tokens en una GPU de $495\cdot 10^{12}$ FLOP/s de pico, con una utilización del 50\,\%. Suponiendo que la pasada hacia delante cuesta $2N$ FLOPs por token y la hacia atrás el doble que la de delante, estime los FLOPs totales y el tiempo de entrenamiento en horas y en días.
\fuente{adaptado de Stanford CS336, primavera de 2026~\cite{cs336a1}, tarea~1, problema \texttt{adamw\_accounting}, apartado~(d).}
\end{ejercicio}

\begin{ejercicio}{3}{eje:4-softmax}
(a) Demuestre la \cref{prop:vjp-softmax} sin usar el jacobiano de la softmax, derivando $\bar z_j=\sum_i\bar p_i\,\partial p_i/\partial z_j$ a partir de $\ln p_i=z_i-\ln\sum_ke^{z_k}$. (b) Deduzca que, para $L=-\ln\softmax(z)_y$, $\partial L/\partial z_i=\softmax(z)_i-\ind[i=y]$, y que las componentes de $\nabla_{\!z}L$ suman~0. (c) Compruébelo con $z=(2,1,0)$ e $y=0$, y dé el valor de $L$.
\end{ejercicio}
